\subsection{Continuity of product measures}

PROPOSITION 5. --- For every measure \(\lambda_0 \in \mathcal{M}(\mathrm{X};\mathbf{C})\) , the mapping \(\mu \mapsto \lambda_0 \otimes \mu\) of \(\mathcal{M}(\mathrm{Y};\mathbf{C})\) into \(\mathcal{M}(\mathrm{X} \times \mathrm{Y};\mathbf{C})\) is vaguely continuous.

For every function \(f \in \mathcal{K}(\mathrm{X} \times \mathrm{Y}; \mathbf{C})\) , we know that the function \(h(y) = \int f(x, y) d\lambda_0(x)\) belongs to \(\mathcal{K}(\mathrm{Y}; \mathbf{C})\) (No.~1, Lemma 2), and \(\langle f, \lambda_0 \otimes \mu \rangle = \langle h, \mu \rangle\) , whence the proposition.

PROPOSITION 6. --- When \(\mathcal{M}(\mathrm{X};\mathbf{C})\) , \(\mathcal{M}(\mathrm{Y};\mathbf{C})\) and \(\mathcal{M}(\mathrm{X}\times\mathrm{Y};\mathbf{C})\) are equipped with the topology of strictly compact convergence (§1, No.~10), the bilinear mapping \((\lambda,\mu)\mapsto\lambda\otimes\mu\) of \(\mathcal{M}(\mathrm{X};\mathbf{C})\times\mathcal{M}(\mathrm{Y};\mathbf{C})\) into \(\mathcal{M}(\mathrm{X}\times\mathrm{Y};\mathbf{C})\) is hypocontinuous for the set of vaguely bounded subsets of \(\mathcal{M}(\mathrm{X};\mathbf{C})\) and \(\mathcal{M}(\mathrm{Y};\mathbf{C})\) (TVS, III, §5, No.~3).

Let \(K \subset X\) , \(L \subset Y\) be two compact sets, A a compact subset of \(\mathcal{K}(X \times Y, K \times L; C)\) , and B a vaguely bounded and closed subset of \(\mathcal{M}(X; C)\) ; it is known that B is vaguely compact (§1, No.~9, Prop. 15), hence also compact for the topology of strictly compact convergence (§1, No.~10, Prop. 17). On the other hand, the Banach space \(\mathcal{K}(X \times Y, K \times L; C)\) is isometric to \(\mathcal{K}(X, K; \mathcal{K}(Y, L; C))\) (No.~1, Lemma 1); the mapping \(\varphi\) of \(\mathcal{K}(X, K; \mathcal{K}(Y, L; C)) \times \mathcal{M}(X; C)\) into \(\mathcal{K}(Y, L; C)\) , such that \(\varphi(g, \lambda)\) is the function h defined by \(h(y) = \int g(x, y) d\lambda(x)\) , is separately continuous by virtue of §3, No.~4, Props. 8 and 9. Since \(\mathcal{K}(X, K; \mathcal{K}(Y, L; C))\) is barreled, it follows that the mapping \(\varphi\) is hypocontinuous relative to the vaguely bounded subsets of \(\mathcal{M}(X; C)\) (TVS, III, §5, No.~3, Prop. 6); the restriction of this mapping to \(A \times B\) is therefore continuous (loc. cit., Prop. 4). The image C of \(A \times B\) under this mapping is consequently a compact subset of the Banach space \(\mathcal{K}(Y, L; C)\) . Now, C is none other than the set of functions \(h(y) = \int f(x, y) d\lambda(x)\) as f runs over A and \(\lambda\) runs over B; by virtue of formula (3) of No.~1, the conditions \(\lambda \in B\) and \(\mu \in C^{\circ}\) therefore imply \(\lambda \otimes \mu \in A^{\circ}\) . In view of the definition of the topology of strictly compact convergence, this proves the proposition (TVS, III, §5, No.~3, Def. 2).

\pandocbounded{\includegraphics[keepaspectratio]{images/part001_f639fdbe.jpg}}

The conclusion of Prop. 6 is no longer valid when the topology of strictly compact convergence is replaced by the vague topology (Exer. 2 c)). However, if B (resp. C) is a vaguely bounded subset of \(\mathcal{M}(\mathrm{X};\mathbf{C})\) (resp. \(\mathcal{M}(\mathrm{Y};\mathbf{C})\) ), then the image of \(\mathrm{B}\times \mathrm{C}\) under the mapping \((\lambda ,\mu)\mapsto \lambda \otimes \mu\) is vaguely bounded in \(\mathcal{M}(\mathrm{X}\times \mathrm{Y};\mathbf{C})\) and therefore the restriction of this mapping to \(\mathrm{B}\times \mathrm{C}\) is vaguely continuous, by virtue of Prop. 6, of §1, No.~10, Prop. 17, and of Prop. 4 of TVS, III, §5, No.~3 (cf.~Exer. 3).

